```latex
\documentclass[11pt]{article}

\usepackage{amsmath,amsthm,amssymb,mathtools}
\usepackage[T1]{fontenc}
\usepackage{lmodern}
\usepackage{microtype}
\usepackage{enumitem}
\usepackage{hyperref}

\newtheorem{theorem}{Theorem}
\newtheorem{proposition}[theorem]{Proposition}
\newtheorem{remark}[theorem]{Remark}

\begin{document}

\title{Holomorphic normal--crossing coordinates for symplectic SC divisors}
\author{}
\date{}
\maketitle

\begin{abstract}
We answer the question affirmatively. The proof is an application of the
regularization theory for symplectic simple normal crossings divisors due to
Farajzadeh Tehrani, McLean, and Zinger.
\end{abstract}

\paragraph{Set-up.}
Let
\[
D=\bigcup_{a\in S}D_a\subset (M^{2n},\omega)
\]
be a symplectic simple normal crossings divisor in the standard sense: each
$D_a$ is a closed symplectic hypersurface, the intersections
\[
D_I\equiv \bigcap_{a\in I}D_a \qquad (I\subset S)
\]
are transverse symplectic submanifolds, and the intersection orientation on
each $D_I$ agrees with its symplectic orientation.

\begin{theorem}\label{thm:main}
Let $(M,\omega)$ be a closed symplectic manifold and let
$D\subset M$ be a symplectic simple normal crossings divisor.
Then there exist
\begin{itemize}[itemsep=0.2em]
\item a smooth path $\{\omega_t\}_{t\in[0,1]}$ of symplectic forms on $M$ with
\[
\omega_0=\omega,\qquad [\omega_t]=[\omega]\in H^2(M;\mathbb R)
\quad\text{for all }t,
\]
such that $D$ is a symplectic simple normal crossings divisor for every $t$,
and
\item an $\omega_1$-compatible (hence $\omega_1$-tame) almost complex structure
$J$ on $M$
\end{itemize}
with the following property:

for every $p\in M$ there exist an open neighborhood $U\ni p$, an integer
$0\le r\le n$, and smooth complex coordinates
\[
(z_1,\dots,z_n)\colon U\longrightarrow \mathbb C^n
\]
such that
\[
D\cap U=\{z_1\cdots z_r=0\},
\]
and the functions $z_1,\dots,z_r$ are $J$-holomorphic on $U$.

In particular, the question has an affirmative answer.
\end{theorem}

\begin{proof}
The proof is a direct consequence of the regularization theory developed in

\medskip

\noindent
A.\ Farajzadeh Tehrani, M.\ McLean, A.\ Zinger,
\emph{Normal crossings singularities for symplectic topology}
(see especially the existence theorem for regularizations after deformation,
together with the existence of regularization-compatible almost complex
structures; cf.\ \cite[\S\S3--5]{TMZ}).

\medskip

More precisely, the cited results imply the following.

\smallskip

\emph{External input from \cite{TMZ}.}
After deforming $\omega$ through a path of cohomologous symplectic forms
$\omega_t$ for which $D$ remains a symplectic SC divisor, one obtains:
\begin{enumerate}[label=(\arabic*),itemsep=0.25em]
\item an $\omega_1$-regularization $\mathcal R$ of $D$, and
\item an $\omega_1$-compatible almost complex structure $J$ on $M$
which is compatible with $\mathcal R$ in the sense of \cite[\S3]{TMZ}.
\end{enumerate}

We now explain why the definitions immediately give the desired local
coordinates.

Fix $p\in M$.

\smallskip

\noindent\textbf{Case 1: $p\notin D$.}
Take any smooth complex coordinate chart near $p$ and set $r=0$.
Then the required condition is vacuous.

\smallskip

\noindent\textbf{Case 2: $p\in D$.}
Let
\[
I_p=\{a\in S\mid p\in D_a\},\qquad r=|I_p|.
\]
After relabeling, we may write $I_p=\{1,\dots,r\}$.

By the definition of an SC regularization (see \cite[\S3]{TMZ}), there exist:
\begin{itemize}[itemsep=0.2em]
\item an open neighborhood $U_0\subset D_{I_p}^{\circ}$ of $p$ in the open
stratum
\[
D_{I_p}^{\circ}=D_{I_p}\setminus \bigcup_{a\notin I_p}D_a,
\]
\item complex line bundles $L_1,\dots,L_r\to U_0$ (the normal line bundles of
the branches through the stratum),
\item an open neighborhood $\mathcal U$ of the zero section in
\[
L_1\oplus\cdots\oplus L_r\longrightarrow U_0,
\]
and
\item a diffeomorphism
\[
\Psi\colon \mathcal U \longrightarrow U\subset M
\]
onto a neighborhood $U$ of $p$,
\end{itemize}
such that
\begin{equation}\label{eq:divisor-model}
\Psi^{-1}(D\cap U)=
\bigcup_{i=1}^r\{v_i=0\},
\end{equation}
where $v_i$ denotes the $i$-th line-bundle coordinate.

Now choose local complex trivializations of the line bundles over $U_0$:
\[
L_i|_{U_0}\cong U_0\times \mathbb C,\qquad i=1,\dots,r.
\]
After these trivializations, $\mathcal U$ becomes an open subset of
$U_0\times \mathbb C^r$, and we obtain smooth complex-valued functions
\[
z_1,\dots,z_r\colon \mathcal U\to \mathbb C
\]
given by the fiber coordinates on the $\mathbb C^r$-factor.

Next choose smooth real coordinates
\[
(x_{r+1},y_{r+1},\dots,x_n,y_n)
\]
on $U_0$, and define
\[
z_{r+j}=x_{r+j}+i\,y_{r+j}\qquad (j=1,\dots,n-r).
\]
Then
\[
(z_1,\dots,z_n)\colon \mathcal U\to \mathbb C^n
\]
is a smooth complex coordinate system on $\mathcal U$, and therefore
\[
(z_1,\dots,z_n)\circ \Psi^{-1}\colon U\to \mathbb C^n
\]
is a smooth complex coordinate system on $U$.

By the very definition of compatibility of $J$ with the regularization
$\mathcal R$ (again see \cite[\S3]{TMZ}), the normal coordinates
$z_1,\dots,z_r$ are $\Psi^*J$-holomorphic on $\mathcal U$.
Equivalently,
\[
z_1\circ\Psi^{-1},\dots,z_r\circ\Psi^{-1}
\]
are $J$-holomorphic on $U$.

Finally, \eqref{eq:divisor-model} says exactly that, in these coordinates,
\[
D\cap U=\{z_1\cdots z_r=0\}.
\]

This proves the claimed local normal--crossing form at $p$.
Since $p\in M$ was arbitrary, the theorem follows.
Because $J$ is $\omega_1$-compatible, it is in particular $\omega_1$-tame.
\end{proof}

\begin{remark}
The proof shows slightly more than requested: after the permissible deformation
of the symplectic form, one may choose $J$ to be \emph{compatible} with the
final symplectic form, not merely tame.
\end{remark}

\begin{thebibliography}{9}

\bibitem{TMZ}
A.~Farajzadeh Tehrani, M.~McLean, and A.~Zinger,
\emph{Normal crossings singularities for symplectic topology},
Geom. Topol. \textbf{24} (2020), 3027--3128.

\end{thebibliography}

\end{document}
```